\documentclass[11pt,a4paper]{article}
\usepackage[T1]{fontenc}
\usepackage[utf8]{inputenc}
\usepackage[english]{babel}
\usepackage{amsmath,amssymb}
\usepackage[margin=2cm]{geometry}
\setlength{\parindent}{0pt}
\setlength{\emergencystretch}{3em}
\allowdisplaybreaks

\title{About \texttt{solution\_sets.tex}}
\author{}
\date{}
\begin{document}
\maketitle
This note explains what the test data in \texttt{solution\_sets.tex} is used for. That file is input for the engine in \texttt{main.c}; it is read by the test \texttt{test\_equations\_inequalities.c} in this folder.

\begin{enumerate}
\item \textbf{Purpose.} The file holds 24 requests of the form $\{x \in \mathbb{R} \mid \ldots\}$: nothing is claimed, the engine must find the set of solutions and write it as a textbook does, for example $x < -3 \lor x \ge 2$.

\item \textbf{What it exercises.} Reading the solutions off the sign table: consecutive rows where the condition holds make an interval or a single point, closed where the endpoint is itself a solution. Irrational endpoints are written with radicals when they are roots of a polynomial of degree two, and get a name otherwise.

\item \textbf{Variety.} No solution, every real number, isolated points, unions of intervals, excluded points, a restricted set such as $\mathbb{R}^{+}$, endpoints $\sqrt{2}$, $1 \pm \sqrt{2}$, $\frac{2 \pm \sqrt{6}}{2}$, a polynomial of degree four whose roots are found by splitting it into two factors of degree two, and roots of cubics, which only get a name. The last six are textbook exercises whose printed result contains a square root: the engine cannot yet confirm that result as a statement (see \texttt{book\_inequalities\_with\_radicals.tex}), but it must find it by itself, written exactly as the book prints it.

\item \textbf{Expected result.} The line \texttt{PASS: 24 of 24 solution sets found and equal to the expected ones}. The expected solutions are written in the test; they were checked by hand or against the textbook, not copied from the engine's output.

\item \textbf{How to run it.} From the top folder of the repository:
\begin{verbatim}
cc -std=c11 -O2 -Wall -Wextra -o /tmp/test_equations_inequalities \
    tests/equations_inequalities/test_equations_inequalities.c \
    && /tmp/test_equations_inequalities
\end{verbatim}
\end{enumerate}
\end{document}
